% context: markscheme (official solutions) for 1-3HW_Measurement
% topic: significant figures, rounding, scientific notation, reading a ruler,
%        unit conversion, average speed
% convention: "=" joins exact equalities; "\approx" introduces a rounded value
\documentclass[12pt, twoside]{article}
\input{../preamble}
\usepackage{float}
\setlength{\parskip}{0.3\baselineskip}

\title{Physics}
\author{Chris Huson}
\date{September 2026}

\fancyhead[LE]{\thepage}
\fancyhead[RO]{\thepage}
\fancyhead[LO]{La Scuola d'Italia / Huson / Physics \\* 18 September 2026}

\begin{document}

\subsubsection*{1.3 Homework Solutions: Measurement quiz review}

\emph{Notation: \,$=$\, joins values that are exactly equal; \,$\approx$\, introduces a
rounded value. Students should write it the same way --- once you round the calculator
display, the equals sign is no longer true. Each conversion runs as one chain of unit
factors and is evaluated once, so there is nothing to round in the middle.}

\begin{enumerate}[itemsep=0.25cm]

\item Significant digits.
  \begin{multicols}{3}
    \begin{enumerate}[itemsep=0.1cm]
      \item \textbf{3}
      \item \textbf{2}
      \item \textbf{3}
      \item \textbf{4}
      \item \textbf{4}
      \item \textbf{4}
    \end{enumerate}
  \end{multicols}
  \emph{(b) The leading zeros in 0.0078 only locate the decimal point; they are not
  measured. (d) The trailing zero in 250.0 is after the decimal point, so it was
  measured --- that is the whole reason for writing it. (e) 0.01070 has leading zeros
  that do not count and a trailing zero that does. (f) Scientific notation settles the
  question for good: everything in the coefficient counts, and nothing else does.}

\item Round to three significant figures.
  \begin{enumerate}[itemsep=0.15cm]
    \item $\sqrt{7} = 2.6457513\ldots \approx \mathbf{2.65}$
    \item $384{,}400~\text{km} \approx \mathbf{384{,}000~\textbf{km}}$, better written
      $\mathbf{3.84 \times 10^{5}~\textbf{km}}$

      \emph{The third figure is the 4 in the thousands place; everything after it goes to
      zero. Writing 384{,}000 is ambiguous --- it looks like it could have six figures ---
      which is exactly why scientific notation is the safer form.}
  \end{enumerate}

\item Scientific notation.
  \begin{enumerate}[itemsep=0.15cm]
    \item $299{,}792{,}458~\text{m/s} = 2.99792458 \times 10^{8}~\text{m/s}
      \approx \mathbf{3.00 \times 10^{8}~\textbf{m/s}}$

      \emph{Write $3.00$, not $3$. The two trailing zeros are the significant figures you
      were asked for.}
    \item $0.000000532~\text{m} = \mathbf{5.32 \times 10^{-7}~\textbf{m}}$
      \quad (already exactly three figures --- nothing to round)
  \end{enumerate}

\item Ordinary decimal numbers.
  \begin{enumerate}[itemsep=0.15cm]
    \item $6.37 \times 10^{6}~\text{m} = \mathbf{6{,}370{,}000~\textbf{m}}$
    \item $6.7 \times 10^{-2}~\text{s} = \mathbf{0.067~\textbf{s}}$
  \end{enumerate}

\item Reading the ruler. The bar runs from 0 to a little past the 7.4 cm mark.
  \begin{enumerate}[itemsep=0.15cm]
    \item $\mathbf{7.45~\textbf{cm}}$ --- the 7.4 is certain (the end lies between the
      4th and 5th millimeter marks past 7), and the final 5 is the estimated digit,
      judged by eye as about halfway between them.

      \emph{Accept 7.44 to 7.46. Do not accept 7.4 (no estimated digit --- it throws away
      information the ruler gives you) or 7.450 (an estimate of an estimate).}
    \item \textbf{Three significant figures.}
    \item $7.45~\text{cm} \times \dfrac{10~\text{mm}}{1~\text{cm}} = \mathbf{74.5~\textbf{mm}}$
      \qquad
      $7.45~\text{cm} \times \dfrac{1~\text{m}}{100~\text{cm}} = \mathbf{0.0745~\textbf{m}}$

      \emph{All three are the same measurement and all three have three significant
      figures. Changing units never changes the precision --- the leading zeros in
      0.0745 are placeholders, not measured digits.}
  \end{enumerate}

\item Convert.
  \begin{enumerate}[itemsep=0.15cm]
    \item $1776~\text{ft} \times \dfrac{1~\text{m}}{3.28~\text{ft}}
      = \dfrac{1776}{3.28}~\text{m} = 541.46\ldots \approx \mathbf{541~\textbf{m}}$

      \emph{The published height is 541.3 m, so this checks. Note that 1776 ft is not a
      measurement at all --- the height was chosen to match the year, so the number is
      exact by design and every digit means something. The 3.28 is the rounded factor
      here, and it is what limits the answer to three figures.}
    \item $40~\text{nickels} \times \dfrac{5.000~\text{g}}{1~\text{nickel}}
      \times \dfrac{1~\text{kg}}{1000~\text{g}} = \mathbf{0.2000~\textbf{kg}}$

      \emph{Keep all four figures: the mint specification is exact to four decimal places
      of a gram and the count of 40 is exact, so nothing here forces rounding. This is why
      a nickel makes a free calibration check for a classroom balance.}
  \end{enumerate}

\item Marathon record: 42.195 km in 2 h 00 min 35 s.

  First the time, in seconds: $2~\text{h} \times \dfrac{3600~\text{s}}{1~\text{h}}
  + 35~\text{s} = 7200~\text{s} + 35~\text{s} = 7235~\text{s}$

  $v = \dfrac{d}{t} = \dfrac{42{,}195~\text{m}}{7235~\text{s}} = 5.8320\ldots
  \approx \mathbf{5.83~\textbf{m/s}}$

  $v = \dfrac{42.195~\text{km}}{7235~\text{s}} \times \dfrac{3600~\text{s}}{1~\text{h}}
  = 20.995\ldots \approx \mathbf{21.0~\textbf{km/h}}$

  \emph{Both start from the same measured distance and time. Converting the 5.83 m/s
  instead is fine too ($\times 3.6$), as long as you carry the unrounded 5.8320 into it.
  Worth saying out loud: 21 km/h is faster than most students can sprint, and the record
  holder keeps it up for two hours.}

\item Frecciarossa at 300 km/h.
  \begin{enumerate}[itemsep=0.15cm]
    \item $300~\dfrac{\text{km}}{\text{h}} \times \dfrac{1000~\text{m}}{1~\text{km}}
      \times \dfrac{1~\text{h}}{3600~\text{s}}
      = \dfrac{300{,}000}{3600}~\dfrac{\text{m}}{\text{s}} = 83.33\ldots
      \approx \mathbf{83.3~\textbf{m/s}}$
    \item $300~\dfrac{\text{km}}{\text{h}} \times \dfrac{1~\text{mi}}{1.61~\text{km}}
      = \dfrac{300}{1.61}~\dfrac{\text{mi}}{\text{h}} = 186.33\ldots
      \approx \mathbf{186~\textbf{mi/h}}$
  \end{enumerate}
  \emph{Rome to Milan is about 570 km and the fast service takes roughly 3 hours, so the
  \emph{average} speed is near 190 km/h --- well under the 300 the train is built for.
  Stops, station approaches, and older track at each end account for the difference. Top
  speed and average speed are different quantities.}

\item Badminton and baseball.
  \begin{enumerate}[itemsep=0.15cm]
    \item Badminton: $493~\dfrac{\text{km}}{\text{h}} \times \dfrac{1000~\text{m}}{1~\text{km}}
      \times \dfrac{1~\text{h}}{3600~\text{s}} = 136.94\ldots \approx \mathbf{137~\textbf{m/s}}$

      Fastball: $95~\dfrac{\text{mi}}{\text{h}} \times \dfrac{5280~\text{ft}}{1~\text{mi}}
      \times \dfrac{1~\text{m}}{3.28~\text{ft}} \times \dfrac{1~\text{h}}{3600~\text{s}}
      = \dfrac{95 \times 5280}{3.28 \times 3600}~\dfrac{\text{m}}{\text{s}}
      = 42.48\ldots \approx \mathbf{42.5~\textbf{m/s}}$
    \item The neat way needs no metric units at all --- the feet and the miles cancel
      each other:

      $t = \dfrac{60.5~\text{ft}}{95~\text{mi/h}} \times \dfrac{1~\text{mi}}{5280~\text{ft}}
      \times \dfrac{3600~\text{s}}{1~\text{h}}
      = \dfrac{60.5 \times 3600}{95 \times 5280}~\text{s} = 0.4342\ldots
      \approx \mathbf{0.43~\textbf{s}}$

      Equivalently, $t = \dfrac{18.44~\text{m}}{42.48~\text{m/s}} \approx 0.43$ s --- same
      answer, more work.

      \textbf{The hitter has about 0.43 s}, and roughly 0.25 s of that is gone before the
      body can respond at all. Less than two tenths of a second remains to judge the pitch
      and start the swing --- and the decision has to be made while the ball is still only
      part way to the plate. This is why hitters talk about reading the pitcher's arm
      rather than watching the ball.
  \end{enumerate}

\item Proton to electron.

  $\dfrac{1.673 \times 10^{-27}~\text{kg}}{9.109 \times 10^{-31}~\text{kg}}
  = \dfrac{1.673}{9.109} \times 10^{-27-(-31)} = 0.18366\ldots \times 10^{4}
  = 1836.6\ldots \approx \mathbf{1837}$

  \textbf{Four significant figures}, because both masses are given to four and a quotient
  carries no more than the least precise value that went into it. The answer is a pure
  number --- the kilograms cancel --- so it has no units.

  \emph{Watch the exponents: subtracting a negative exponent adds. Students who write
  $10^{-58}$ have multiplied instead of divided; a quick sanity check is that the proton is
  obviously the heavier particle, so the answer must come out well above 1.}

\item The Atlantic, widening 2.5 cm/yr.
  \begin{enumerate}[itemsep=0.15cm]
    \item $2.5~\dfrac{\text{cm}}{\text{yr}} \times \dfrac{1~\text{m}}{100~\text{cm}}
      \times \dfrac{1~\text{yr}}{3.156 \times 10^{7}~\text{s}}
      = \dfrac{2.5}{100 \times 3.156 \times 10^{7}}~\dfrac{\text{m}}{\text{s}}
      = 7.921\ldots \times 10^{-10} \approx \mathbf{7.9 \times 10^{-10}~\textbf{m/s}}$

      \emph{Two significant figures --- the 2.5 has two. This is a number that only makes
      sense in scientific notation: written out it is 0.00000000079 m/s, and nobody can
      read that.}
    \item $1.0~\text{m} \times \dfrac{100~\text{cm}}{1~\text{m}}
      \times \dfrac{1~\text{yr}}{2.5~\text{cm}} = \mathbf{40~\textbf{yr}}$

      \emph{So the Atlantic has widened by about a meter since these students' parents
      were born --- a rate that is imperceptible per second and obvious per lifetime.}
  \end{enumerate}

\item Three measurements of the hallway.
  \begin{enumerate}[itemsep=0.15cm]
    \item $\dfrac{24.1 + 24.0 + 24.2}{3}~\text{m} = \dfrac{72.3}{3}~\text{m}
      = \mathbf{24.1~\textbf{m}}$
    \item The three readings disagree in the tenths place, so the tenths place is already
      the uncertain digit --- the hallway is $24.1 \pm 0.1$ m. Writing 24.1000 m claims the
      length is known to a tenth of a millimeter, which is a thousand times better than the
      measurements actually agree with each other.

      \emph{The extra zeros came from dividing by 3 on a calculator, not from the tape
      measure. Averaging repeated measurements does make the result more reliable, but it
      cannot invent precision the instrument never had. The correct report is 24.1 m ---
      three significant figures, the same as each reading.}
  \end{enumerate}

\item \textbf{Challenge: gasoline.}
  \begin{enumerate}[itemsep=0.15cm]
    \item Two conversions in one chain --- currency and volume:

      $1.82~\dfrac{\text{\texteuro}}{\text{L}}
      \times \dfrac{\$1.09}{\text{\texteuro}1}
      \times \dfrac{3.785~\text{L}}{1~\text{gal}}
      = 1.82 \times 1.09 \times 3.785~\dfrac{\$}{\text{gal}} = 7.5086\ldots
      \approx \mathbf{\$7.51~\textbf{per gallon}}$

      \emph{Set the factors up so the euros cancel and the liters cancel, leaving dollars
      over gallons. If the answer comes out near \$0.44, one of the two factors was
      inverted --- and the units in the cancelling will show which.}
    \item $\$7.51 - \$3.45 = \mathbf{\$4.06~\textbf{more per gallon}}$, roughly
      \textbf{2.2 times} the New York price.

      \emph{Most of that gap is tax, not the cost of the fuel.}
    \item $45~\text{L} \times \dfrac{\text{\texteuro}1.82}{1~\text{L}}
      \times \dfrac{\$1.09}{\text{\texteuro}1} = \mathbf{\$89.27}$

      \emph{Or \texteuro81.90 before converting. Money is the one place where rounding to
      the cent is expected regardless of significant figures --- a price is an exact
      amount, not a measurement.}
  \end{enumerate}

\end{enumerate}
\end{document}
